\chapter{Introduction}
\label{ch1}

\graphicspath{{Chapter1/}}

In software engineering, formal modeling notations occupy a fundamental niche due to their importance in the early phases of requirements definition, where they are used to represent and convey to developers and stakeholders essential information such as conceptual models, use cases, goals, and formal requirements, design and architecture pattern, process execution, or infrastructure and deployment choices. Despite the importance of modeling languages in the software development process, they tend to be underused in professional practice, with developers preferring more informal methods to communicate requirements, and their importance is not conveyed successfully in educational curricula, with modeling topics being perceived as frustrating and inaccessible by novice learners.

Furthermore, teaching modeling notations such as the Unified Modeling Language (UML) or Business Process Model and Notation (BPMN) implies that students have to perform practical assignments that need continuous feedback loops that can identify recurring error patterns and provide guidance on how to avoid them by pointing out correct modeling practices; however, tools that can support educators by providing large-scale feedback are limited. This thesis aims at addressing this gap by exploring whether the use of gamification can be used as a supporting technique to enhance the learning process focused on modeling notations, with the aim of filling the research gap related to reduced body of rigorously evaluated gameful solutions for software modeling. A third gap identified at the beginning of the thesis work was the relatively unexplored potential of large language models as assistants for both learners and educators in diagram evaluation and generation. The work described in this thesis covers the intersection of these three gaps by asking whether gamified modeling tools can improve student engagement, performance, and learning progression, as well as asking whether intelligent language model agents can complement or extend the automated evaluation mechanisms that are necessary for gamified tools to function. 

The challenges in learning software modeling notations' fundamentals, together with the obstacles that make acquiring proficiency a daunting task for students, are well documented by the research literature: surveys on students and practitioners that study and use UML identified multiple issues that make it hard to justify adopting the notation such as its inherent complexity, the need to put high effort to ensure diagrams remain consistent following code evolutions, and a lack of valid interactive learning tools \cite{siau2006identifying}. Other investigations, focused explicitly on student behavior identified systematic error patterns that emerge as a consequence of the reluctance to learn, patterns that relate to incorrect understanding of association and inheritance semantics, inability to correctly represent object identity, and incorrect placement of attributes; these errors tend to persist even following explicit teacher guidance and impact the overall quality of the resulting student-produced models \cite{boustedt2012students}. Other obstacles toward adopting UML in agile context include the preference for informal, easier to maintain representations, preferring a focus on speed over precision, excessive tool complexity discouraging continuous use, and low return on investment \cite{stettina2012documentation, oliveirajr2025uml}.

For what concerns BPMN, obstacles to adoption are comparable: novice modelers tend to struggle with correctly representing control flow executions, understanding the semantics behind gateways, and tend to have trouble handling communication with multiple pools; these challenges can also be further amplified by the lack of standardized pedagogical frameworks and by the anxiety-inducing assignments tied to the aforementioned obstacles that hinder formal progress representation \cite{bandara2010bpm, maslov2022towards}. All these obstacles share a common structural challenge: reaching an adequate level of competence in software modeling tasks requires sustained practice supported by immediate, specific, and actionable feedback that can help students identify the causes behind their errors and provide insights on how to correct them instead of simply reporting failure \cite{bogdanova2019towards}. However, providing feedback with such a level of detail in courses with hundreds of enrolled students is not feasible through instructor assistance alone: automated and interactive feedback can be considered strategies with high potential for improving learning outcomes in modeling education, with the caveat that the quality and immediacy of the feedback must be seen as the primary factor that affects student behavior \cite{huber2022tool}.

Gamification, formally defined as the act of using game elements in non-recreational contexts for the purpose of improving performance and increasing motivation and participation \cite{deterding_game_2011}, is a technique that has been proposed in response to the motivational challenges behind the low student interest in modeling activities. There are multiple examples in the literature that motivate why a careful selection of game mechanics could be beneficial: elements such as progress bars and experience points support the students' need for competence and make skill development tangible, avatar customization and narrative elements allow for autonomous experiences and make students feel connected with their tasks, while competition and cooperation mechanics foster the need for relatedness \cite{deci2000intrinsic}. 

Evidence from studies on gamification applied to education suggests that the technique can bring small to medium improvements on cognitive outcomes and engagement when the game mechanics are selected to match the educational process, but also highlights the risk of poorly designed implementations that focus excessively on extrinsic rewards or competitive comparison, with consequent loss of motivation over time and neutral, or even negative results \cite{hamari2014does, sailer2020gamification, hanus2015assessing}. In the software engineering field, studies that apply gamification focus most on coding, development, and testing, while design, modeling, and requirements engineering tasks are substantially less explored \cite{pedreira2015gamification}.

This gap is particularly significant in terms of tools designed to support BPMN and UML modeling: only one gamified tool exists for BPMN modeling, and its intended use is for industrial modeling with a focus on sustainability-oriented assessment rather than educational use \cite{mancebo_bpms-game_2017}, while most tools that apply gamification to UML tend to cover class diagrams, are serious games (games with an educational purpose rather than educational tools with game elements), and make use of predefined element selection instead of allowing free diagram construction; furthermore, at the time the PhD research began, almost none of them were subjected to longitudinal controlled evaluations. The absence of rigorously evaluated gamified educational environments that provide automated formative feedback constitutes the main research gap that motivated the first research goal of the thesis.

The second research goal, instead, emerged from the limitations found with the empirical validations performed to address the first goal: the tools developed during the PhD implemented evaluation systems that would assess a student's diagram to identify content matching the intended correct solutions well enough to receive rewards and obtain corrective feedback for errors such as missing content or imperfections related to the exercise's target domain. The matching procedures implemented within these evaluation engines relied on matching the labels assigned to diagram elements against teacher-defined reference structures using normalized edit distance comparisons to identify content intended to represent specific requirements, with synonyms and alternative spelling options supported to accommodate naming variations. This approach proved to be effective in most scenarios but encountered issues in cases where a student would label a given element using a semantically valid phrase/noun that was, however, not expected in the defined reference, leading to false negatives (e.g., elements correctly modeled marked as missing), resulting in unjustified errors that would cause frustration and reduce the educational impact of the feedback mechanism. 

Large language models (LLMs), a technology that became available shortly after the beginning of the PhD, showed promising capabilities in semantic similarity assessment thanks to their massive linguistic training, and were identified as a potential solution to the excessively rigid evaluation systems. Furthermore, LLMs also raised a separate empirical question: could a general-purpose model, with its text interpretation capabilities, be used as an educational assistant and generate diagrams of sufficiently high quality to serve as draft example solutions, exercise templates, or illustrations of domain models to be shown to students?

Early surveys of LLM use in software engineering tasks confirmed that researchers prefer focusing on code centric activities, while model-based tasks were relatively less popular \cite{fan2023large, hou2023large}, and the early studies on LLM-generated UML models pointed out a limitation: the produced diagrams were structurally sound and adhered to the UML language's syntactic rules, while also presenting significant structural and semantic deficiencies such as incorrect association types and multiplicities, requiring constant human correction \cite{camara2023assessment, chen2023automated}. These gaps motivated the second research goal of the PhD, and led to the studies described within Chapter \ref{ch6}.

The research activities described in this thesis evolved progressively across multiple phases, with each one motivated by the limitations that had emerged in the preceding ones. The first phase saw the development of a prototype web application for gamified BPMN exercise solving named BIPMIN; the prototype was designed with three distinct themes guiding the selection of game mechanics, in order to identify the most effective motivational approaches in a modeling context. The prototype was first tested with a preliminary usability study with twelve participants \cite{bedwell_bipmin_2023}, and a revised version combining the most appreciated mechanics was evaluated with a controlled experiment in January 2023. The study involved Masters-level students and aimed at evaluating how gamification would impact diagram correctness, student productivity, and user perception \cite{garaccione_gamification_2023, garaccione_gamification_2024}; findings from the study confirmed a positive effect on semantic diagram quality and involvement with the exercises, but also exposed a second critical limitation, in addition to the aforementioned rigidity of the evaluation engine which caused excessive and unfair penalization for semantically correct modeling choices: a gamified tool that implements long-term mechanics such as leaderboards or level progression systems cannot be fully evaluated with a single-session design, since those mechanics cannot be completely appreciated, and their impact on student performance cannot be observed.

These two limitations guided the second phase, where a generalized gamification framework was designed, distinguishing between exercise-specific mechanics that assist students in the context of each separate exercise and long-term mechanics that motivate students to keep interacting with the gamified tool and, in turn, perform more practical activities. A new tool, UMLegend \cite{cagnazzo_2023_umlegend}, was developed according to this new framework focused on UML class diagram evaluation and with a tailor-made, less restrictive, evaluation engine supporting multiple reference solutions and additional naming variability: this tool was then evaluated with a controlled experiment performed in January 2024 \cite{garaccione_2024_sqj, garaccione_2026_gamification}.

In parallel with the development of this new tool, the BIPMIN application was also redeveloped from scratch following the tenets of the new framework and then deployed longitudinally throughout the 2023-2024 edition of an Information Systems course, providing the first example of longitudinal evaluation of gamified process modeling, deriving insights on whether the technique sustains positive effects over an entire academic semester and impacts final evaluation grades \cite{garaccione2024ease, garaccione2024esem}.

In both controlled experiments on the two tools, the limitation of strict evaluation engines was still observed, although with less frequency compared to the first empirical study and reduced impact on student frustration; the need to include element matching approaches within the tools, together with the emergence of large language models, motivated the third phase of the research: a series of three empirical studies examining LLM agents' capabilities to generate UML diagrams, together with the development of an intelligent, LLM-based assistant for the creation of enhanced reference solutions to provide additional variability for students solving exercises using the UMLegend tool \cite{garaccione2026automated}. The first two empirical studies, started as Masters' thesis works in late 2023, instead covered the comparison of LLM-generated class diagrams against student-produced ones \cite{debari2024evaluating} and an analogous evaluation focused on use case diagrams \cite{garaccione2025vega}; the findings that emerged from those studies pointed out that LLMs could generate syntactically valid diagrams, although with a tendency to hallucinate superfluous elements and with critical semantic violations that made them incomplete. The final study described in this work was also conceptualized as a Master thesis at first, and aimed at overcoming a potential limitation of the two aforementioned studies: the use of unstructured minimal prompts may have impacted the quality of the produced models, and advanced techniques could let LLM agents create more suitable diagrams. For this purpose, a multi-model comparative study was performed in 2025, adopting role-based few-shot prompting and using four LLM architectures, covering both proprietary and open source models to identify each agent's strengths and weaknesses \cite{garaccione2025models}.

The research is organized around six overarching research questions that span across the three research phases and are formulated at an abstraction level that encompasses the chapter-specific research questions, with the goal to provide a unified evaluative flow across the entire thesis. Three questions concern the gamification efforts conducted in the first two years of the PhD:

\begin{itemize}
    \item[\textbf{RQ1:}] \textit{Does the application of gamification to software modeling exercises lead to measurable improvements in student performance?}\\
    This question is answered with the empirical studies described in Chapters \ref{ch3} and \ref{ch4}: the use of within-subjects crossover experiments that compare gamified and non-gamified modeling assignments allows to compare student performance and makes it possible to identify differences in the outcomes tied to the treatment. The studies aim to measure whether gamification affects the productivity aspects that lead students to engage with the modeling environments, as well as the impact on diagram quality, measured with completeness, and syntactic and semantic quality. In the BIPMIN study described in Chapter \ref{ch3}, semantic quality was measured through diagram completeness, while subsequent studies measured it through counts of domain representation errors; the two measures are directionally opposite and not directly comparable.
    \item[\textbf{RQ2:}] \textit{Does sustained exposure to a gamified modeling tool over the duration of an academic course produce improvements in modeling competence?}\\
    This question, reliant on the longitudinal study presented in Chapter \ref{ch5}, moves the temporal dimension beyond single interactions with gamified tools, and measures the impact brought by gamification applied in a longitudinal setting on the grades obtained by students exposed to the tool in final course evaluations through a comparison with students of the preceding edition of the course acting as the non-gamified baseline. Moreover, the longitudinal study measures the impact gamification has on skill growth through comparisons of student performance over the set of exercises offered throughout the course.
    \item[\textbf{RQ3:}] \textit{How do students perceive gamified modeling tools?}\\
    This question spans across all three controlled studies and is answered with the collection and analysis of student opinions through validated questionnaires that measure usability, gamification-specific perception, appreciation of the implemented game mechanics and their perceived impact, and design improvements and suggestions to improve the tool for future use.
\end{itemize}

The second set of research questions is instead structured around the studies on LLM-assisted modeling activities, and is presented below:

\begin{itemize}
    \item[\textbf{RQ4:}] \textit{Can LLM agents generate UML class and use case diagrams whose quality is comparable to that of student-produced diagrams?}\\
    This question is evaluated through two comparative studies where automatically generated diagrams are compared against solutions drawn by students working under exam-like conditions, with assessment criteria focused around syntactic correctness, semantic adherence to the target domain, closeness to expert approved solutions, and presence of hallucinations expressed as superfluous elements. This question is of pivotal importance for the third research phase, since it establishes whether the quality of LLM-generated artifacts is acceptable enough to integrate the technology into the educational workflow, or whether the differences in quality are too large to be surpassed without extensive human correction.
    \item[\textbf{RQ5:}] \textit{Can an LLM-based tool reduce the effort required by educators to create exercise reference solutions?}\\
    This question guided the development of a prototypal assistant that integrates LLM support in the process for defining reference solutions in UMLegend, with a proof of concept evaluation that observes whether the generated draft solutions are structurally correct and can be used to support an interactive workflow between human lecturers and LLMs, potentially more efficient than completely manual generation. Given the effort needed to define solutions and the issues that come with under-specified references, the development of such a tool is crucial to extend the capabilities of the tools developed in the first two research phases.
    \item[\textbf{RQ6:}] \textit{How do different LLM architectures compare in diagram generation quality under equal conditions (exercise selection and prompting strategy)?}\\
    This question builds upon the findings of the first two comparative studies on LLM-assisted diagram generation, introducing a more advanced prompting strategy tested on four different LLM agents, with the intent of identifying the best suited models to integrate in tool development and educational activities depending on the specific model-related goal.
\end{itemize}

The contributions to the research area of software engineering education made to answer these research questions can be summarized according to four dimensions. The empirical contribution consists of seven controlled studies designed using the Goal-Question-Metric template, ensuring a traceable alignment between each study's goals, the research questions formulated to reach them, the metrics used to answer the questions, and the statistical analyses on said metrics. Validated instruments were also applied throughout the controlled studies to measure participant perception such as the Technology Acceptance Model, the GAMEX scale to measure the perception of the gamified experience, and the System Usability Scale; the consistent use of the same instruments allows for meaningful cross-experiment comparison \cite{basili1994gqm}, necessary to draw findings and answer the research questions.

The second contribution, more focused on tool development, introduces two developed and empirically validated software artifacts in the form of the BIPMIN and UMLegend gamified platform, with the former providing a dual contribution, first as a prototype and then as a redesigned longitudinal application; the prototypal intelligent assistant, named TIGRE, also falls within this category.

The design contribution is finalized through the articulation of a gamification framework that distinguishes exercise-specific feedback mechanics, whose main purpose is providing actionable suggestions for students to correct their mistakes and improve their knowledge, and mechanics focused on long-term engagement that motivate students to continue using educational tools. The validity of the framework and of the design choices that led to its formulation are validated by the consistent findings across the three empirical studies, where the different categories of game mechanics produced different motivational effects depending on the context.

Lastly, the methodological contribution consists of the application of three-dimensional quality assessment criteria \cite{lindland1994understanding} across all studies: the distinction between syntactic, semantic, and pragmatic quality is followed in both the evaluation engines used in the gamified tools and in the manual evaluation protocols devised for the LLM studies, enabling a consistent comparison of human and agent modeling performance across different studies and different notations.

The thesis is organized as follows: Chapter \ref{ch2} provides the theoretical background on which the research activities are built upon by introducing Model-Driven Engineering, with a focus on the notations explored by the subsequent studies, then presents gamification theory, its applications in the educational and software engineering domains, followed by a mapping of existing strategies focused on conceptual modeling education; it then introduces large language models and their applications to software modeling activities, before presenting the shared empirical methods used across the studies.

Chapter \ref{ch3} presents the first iteration of the BIPMIN prototype, describing the preliminary usability study, the redesign process performed before the 2023 empirical study, and summarizes the results achieved by said experiment, with the lessons learned that guided the subsequent research activities. The framework designed as a consequence of these findings is described in Chapter \ref{ch4}, which discusses how it was implemented in the new UMLegend tool and then follows with the description of the controlled experiment performed with the new tool, extending the body of contributions brought by the thesis.

Chapter \ref{ch5} then continues with the description of the improved BIPMIN tool and its longitudinal application in the 2023-2024 academic year, presenting the results of the empirical study and confirming findings related to the gamification aspect of the PhD work. The LLM-focused research activities are instead presented in Chapter \ref{ch6}, that begins from the first feasibility evaluation to the development of the intelligent solution creator before delving into the multi-model assessment. 

Lastly, Chapter \ref{ch7} consolidates the lessons learned across the entire set of research work, draws conclusions, and describes the planned future research activities intended to extend the thesis work to contribute to the model-based engineering research community.
